\documentclass[10pt,a4paper]{amsart}
\usepackage[margin=19mm]{geometry}
\usepackage{amsmath,amssymb,mathtools}
\usepackage[scr=rsfs,cal=euler]{mathalfa}
\usepackage{xfrac}
\usepackage[unicode,hidelinks,bookmarks=false]{hyperref}
\newcommand{\ie}{i.e.\ }
\newcommand{\cf}{cf.\ }
\newcommand{\resp}{respectively\ }
\newcommand{\Subsection}{\subsection}
\providecommand{\nobreakdash}{\nobreak\hbox{-}}
\setlength{\emergencystretch}{2em}
\multlinegap0pt
\usepackage{bm}
\usepackage{amssymb}
\usepackage{mathtools}
\usepackage[shortlabels]{enumitem}
\newtheorem{theorem}{Theorem}
\newtheorem{lemma}{Lemma}
\newcommand{\AND}{\quad{\rm and}\quad}
\newcommand{\E}{\mathbb{E}}
\newcommand{\Gam}{\Gamma}
\newcommand{\Jfour}{\mathcal J_4}
\newcommand{\NN}{\mathcal{N}}
\newcommand{\R}{\mathbb{R}}
\newcommand{\cC}{\mathcal{C}}
\newcommand{\cZ}{\mathcal{Z}}
\newcommand{\dW}{d_{\rm W}}
\newcommand{\dd}{d}
\newcommand{\eps}{\varepsilon}
\newcommand{\ind}{\mathbf{1}}
\newcommand{\noi}{\noindent}
\title{Gamma approximation and Poisson--Gaussian invariance principle on Poisson chaos}
\author{Dionysis Milesis, Guangqu Zheng}
\date{Source: arXiv:2609.02136v1 (CC BY 4.0)}
\begin{document}
\maketitle

\noindent\emph{Source and licence.} Gamma approximation and Poisson--Gaussian invariance principle on Poisson chaos. Version arXiv:2609.02136v1 (CC BY 4.0), distributed by arXiv under the Creative Commons Attribution 4.0 International licence, http://creativecommons.org/licenses/by/4.0/, as recorded in the arXiv licence metadata for this exact version and shown on the abstract page https://arxiv.org/abs/2609.02136v1. Full licence notice: \texttt{licenses/arxiv-2609.02136-CC-BY-4.0.txt}.\par\smallskip

\noindent\textit{Editorial scope.} Counted unit: the article's theorem thm4 (source0.tex lines 625--671). The article's complete original proof of it is delivered with it (source0.tex lines 2701--2904, one \texttt{proof} environment of 204 lines, which the article places away from the statement). No mathematics is written, repaired or re-derived: the statement and the proof are the article's own lines, in the article's own notation, and no retained line is substituted or re-flowed.  Retained uncounted as the source-local context the proof needs: lines 320--324; lines 1041--1047; lines 1512--1516; lines 1647--1653; lines 2521--2525; lines 2547--2565; lines 2664--2677.  Nothing was omitted from any retained span, so the package carries no omission record.  No retained line contains a citation, so the wrapper carries no bibliography and no bibliography entry.  No retained line contains a figure environment, an image-inclusion command or drawing code, so no image is delivered and the package declares no asset dependency.  Editorial formatting only, in addition: an AMS \texttt{amsart} preamble that restates the article's own declarations and macros used by the retained lines, namely the environments \texttt{theorem}, \texttt{lemma} on separate counters, together with the standard packages those lines need.  The retained environments print in this document as Theorem 1, Lemma 1 in order of appearance, whereas the article prints the same items with the numbering of its own full document.  The complete native source of the article as distributed in the arXiv e-print for this exact version is bundled unchanged as \texttt{sources/209177/source0.tex}.
\begin{align} \label{J4F}
\Jfour(F)
=
\int_\cZ\E[|D_zF|^4]\mu(\dd z).
\end{align}

\begin{align}
 \int_{\cZ^r}\E[|D^r_{\boldsymbol z}F|^2]
 \,\mu^r(\dd\boldsymbol z)
 =\frac{q!}{(q-r)!}\E[F^2],
 \quad 1\leq r\leq q.
 \label{D_L2}
\end{align}

\begin{align}\label{def_dW}
 \dW(X,Y)
 =\sup_{\|h'\|_\infty\leq1}
 |\E[h(X)]-\E[h(Y)]|.
\end{align}

\begin{align}
 \Gam(F_n,F_n)&\longrightarrow\Gam(F,F)
 \quad\text{in }L^2(\Omega),
 \label{core_Gamma}\\
 \Jfour(F_n)&\longrightarrow\Jfour(F).
 \label{core_J4}
\end{align}

\begin{align}\label{FG_1}
 F=I_q^\eta(f)\in\cC_q^\eta\cap L^4(\Omega)
 \AND
 G=I_q^W(f)\in\cC_q^W.
\end{align}

\begin{lemma}
\label{lem_tetra_core}
Let $F$ and $G$ be as in \eqref{FG_1}.  After passing,
if necessary, to the marked control space
$(\cZ\times[0,1],\mu\otimes\dd u)$, there are bounded symmetric
tetrahedral step kernels $f_m$ of finite-measure support such that,
with $F_m=I_q^\eta(f_m)$ and $G_m=I_q^W(f_m)$,
\noi
\begin{align}
 \|F_m-F\|_{L^4(\Omega)}
 +\|G_m-G\|_{L^4(\Omega)}
 &\longrightarrow0,
 \label{tetra_L4}\\
 \int_\cZ\E[|D_z(F_m-F)|^4]\mu(\dd z)
 &\longrightarrow0.
 \label{tetra_D4}
\end{align}
The partition cells of the kernels $f_m$ are chosen to have measure at most one.
\end{lemma}

\begin{lemma}
\label{lem_hyb_4}
Let $R$ be a multilinear polynomial of degree at most $r$ in mutually
independent centered coordinates, each of which is either Gaussian or
centered Poisson.  On a product extension, replace every Gaussian
coordinate by a centered Poisson coordinate with the same variance,
independent of all original coordinates and of the other replacement
coordinates.  Denote the resulting polynomial by $R^\eta$.  Then
\noi
\begin{align}
 \E[|R|^4]\leq3^{2r}\E[|R^{\eta}|^4].
 \label{hyb_4}
\end{align}
\end{lemma}

\begin{theorem}
\label{thm4}
Fix an integer $q\geq1$.  Let $f\in L_s^2(\mu^q)$ and set
\noi
\begin{align*}
 F=I_q^\eta(f)\in\cC_q^\eta\cap L^4(\Omega),
 \quad
 G=I_q^W(f)\in\cC_q^W,
 \AND
 \sigma^2=\E[F^2]=\E[G^2]
 =q!\|f\|_{L^2(\mu^q)}^2.
\end{align*}
There exists a constant $K_q<\infty$, depending only on $q$, such
that, for every $h\in C^3(\R)$ with bounded derivatives up to order
three,
\noi
\begin{align}
 \left|\E[h(F)]-\E[h(G)]\right|
 \leq
 K_q\|h^{(3)}\|_\infty
 \sqrt{\sigma^2\Jfour(F)}.
 \label{thm5_1}
\end{align}
Moreover, there exists a constant $L_q<\infty$, depending only on
$q$, such that
\noi
\begin{align}
 \dW(F,G)
 \leq
 L_q\bigl\{\sigma^2\Jfour(F)\bigr\}^{1/6},
 \label{thm5_2}
\end{align}
where $\dW$ denotes the $1$-Wasserstein distance as in \eqref{def_dW}.
Consequently, if $F_n=I_q^\eta(f_n)$
and
 $G_n=I_q^W(f_n)$
satisfy
\noi
\begin{align} 
 \sup_n\E[F_n^2]<\infty
 \AND
 \Jfour(F_n)\longrightarrow0,
  \notag%\label{thm5_3}
\end{align}
then $\dW(F_n,G_n)\longrightarrow0$.

\end{theorem}

\begin{proof}[Proof of Theorem~\ref{thm4}]

We first prove \eqref{thm5_1}.  By Lemma~\ref{lem_tetra_core},
\eqref{tetra_L4}, \eqref{tetra_D4}, and the reverse triangle
inequality for $\Jfour^{1/4}$ used in \eqref{core_J4}, it suffices to
consider a tetrahedral step kernel relative to pairwise disjoint cells
$A_1,\ldots,A_N$ satisfying $\lambda_i:=\mu(A_i)\leq1$.
Set

\noi
\begin{align*}
 X_i:=\eta(A_i)-\lambda_i
 \AND
 Y_i:=W(A_i).
\end{align*}
Then the two families $(X_i)_{i=1}^N$ and $(Y_i)_{i=1}^N$ are
mutually independent, with $\E[X_i]=\E[Y_i]=0$ and
$\E[X_i^2]=\E[Y_i^2]=\lambda_i$.
Since the kernel is tetrahedral, there is a homogeneous multilinear
polynomial $Q$ of degree $q$ such that
$ F=Q(X_1,\ldots,X_N)$ and $G=Q(Y_1,\ldots,Y_N)$.


Next, we use a Lindeberg replacement argument.  For $i=0,\ldots,N$, define
the hybrid vector
\noi
\begin{align*}
 Z^{(i)}=(Y_1,\ldots,Y_i,X_{i+1},\ldots,X_N).
\end{align*}
Thus, $Z^{(0)}=(X_1,\ldots,X_N)$ and
$Z^{(N)}=(Y_1,\ldots,Y_N)$, and therefore
\noi
\begin{align}
 \E[h(F)]-\E[h(G)]
 =
 \sum_{i=1}^N
 \left\{
 \E[h(Q(Z^{(i-1)}))]
 -
 \E[h(Q(Z^{(i)}))]
 \right\}.
 \label{Lind_tel}
\end{align}

Since $Q$ is multilinear, for each $i$ we may write
\noi
\begin{align*}
 Q(x_1,\ldots,x_N)
 =
 U_i(x_1,\ldots,\widehat{x_i},\ldots,x_N)
 +
 x_iV_i(x_1,\ldots,\widehat{x_i},\ldots,x_N),
\end{align*}
where $U_i$ and $V_i$ do not depend on the $i$th coordinate.
Evaluate $U_i$ and $V_i$ at
\noi
\begin{align*}
 (Y_1,\ldots,Y_{i-1},X_{i+1},\ldots,X_N)
\end{align*}
and denote the resulting random variables by
$U_i^{\rm hyb}$ and $V_i^{\rm hyb}$.  We may therefore write

\noi
\begin{align*}
 Q(Z^{(i-1)})
 =
 U_i^{\rm hyb}+X_iV_i^{\rm hyb}
 \AND
 Q(Z^{(i)})
 =
 U_i^{\rm hyb}+Y_iV_i^{\rm hyb}.
\end{align*}
Moreover, $(U_i^{\rm hyb},V_i^{\rm hyb})$ is independent of both
$X_i$ and $Y_i$.  Taylor's formula gives, for every $x\in\R$,

\noi
\begin{align*}
 h(U_i^{\rm hyb}+xV_i^{\rm hyb})
 &=
 h(U_i^{\rm hyb})
 +h'(U_i^{\rm hyb})xV_i^{\rm hyb}
 +\frac12h''(U_i^{\rm hyb})x^2
 (V_i^{\rm hyb})^2
 +R_i(x),
\end{align*}
where
$|R_i(x)|\leq\frac16\|h^{(3)}\|_\infty
|x|^3|V_i^{\rm hyb}|^3$.
Since $X_i$ and $Y_i$ have matching first two moments, the constant,
linear, and quadratic terms cancel after taking expectations
so that

\noi
\begin{align}
 &\left|
 \E h(U_i^{\rm hyb}+X_iV_i^{\rm hyb})
 -
 \E h(U_i^{\rm hyb}+Y_iV_i^{\rm hyb})
 \right|
\leq
 \frac{1}{6} \| h^{(3)}\|_\infty
 \E[|V_i^{\rm hyb}|^3]
 \left(
 \E|X_i|^3+\E|Y_i|^3
 \right).
 \label{Lind_step}
\end{align}


A direct computation, using $\lambda_i\leq1$, gives

\noi
 \begin{align} \label{coord_third}
\max\{\E[|X_i|^3],\E[|Y_i|^3]\}\leq2\lambda_i.
\end{align}


Let $V_i^\eta$ be the polynomial $V_i$ evaluated at the all-Poisson
vector obtained from $(X_1,\ldots,X_N)$ by deleting the $i$th
coordinate.  Since $V_i$ has degree at most $q-1$, Lemma~\ref{lem_hyb_4}
and equality of the coordinate variances give
\noi
\begin{align}
 \E[(V_i^{\rm hyb})^4]
 &\leq3^{2(q-1)}\E[(V_i^\eta)^4],
 \quad
 \E[(V_i^{\rm hyb})^2]=\E[(V_i^\eta)^2].
 \label{Vi_hybrid}
\end{align}


Since $F=Q(X_1,\ldots,X_N)$ is multilinear,
\noi
\begin{align*}
 D_zF=\sum_{i=1}^N\ind_{A_i}(z)V_i^\eta.
\end{align*}
We deduce from \eqref{D_L2} with $r=1$ and from the definition
\eqref{J4F}, respectively, that
\noi
\begin{align}
 \sum_{i=1}^N\lambda_i\E[(V_i^\eta)^2]
 =q\sigma^2
 \AND
 \sum_{i=1}^N\lambda_i\E[(V_i^\eta)^4]
 =\Jfour(F).
 \label{infl_fourth}
\end{align}

By Cauchy--Schwarz, \eqref{Vi_hybrid}, and \eqref{infl_fourth},
we obtain

\noi
\begin{align*}
 \sum_{i=1}^N
 \lambda_i\E[|V_i^{\rm hyb}|^3]
 &\leq
 \left(
 \sum_{i=1}^N
 \lambda_i\E[(V_i^{\rm hyb})^2]
 \right)^{1/2}
 \left(
 \sum_{i=1}^N
 \lambda_i\E[(V_i^{\rm hyb})^4]
 \right)^{1/2}\\
 &\leq
 3^{q-1}
 \sqrt{q\sigma^2\Jfour(F)}.
\end{align*}
Combining this with
\eqref{Lind_tel}, \eqref{Lind_step}, and
\eqref{coord_third}, we obtain
\noi
\begin{align*}
 |\E[h(F)]-\E[h(G)]|
 &\leq
 3^{q-1}\sqrt q\,
 \|h^{(3)}\|_\infty
 \sqrt{\sigma^2\Jfour(F)}.
\end{align*}
This proves \eqref{thm5_1} with $K_q=3^{q-1}\sqrt q$.


\medskip

We next prove \eqref{thm5_2}.  Let $N\sim\NN(0,1)$ and, for a
one-Lipschitz function $h$, define
$h_\eps(x)=\E[h(x+\eps N)]$.  Differentiating under the expectation
and using $\E|N|\leq1$ and $\E|N^2-1|\leq2$, we obtain
$\|h-h_\eps\|_\infty\leq\eps$ and
$\|h_\eps^{(3)}\|_\infty\leq2\eps^{-2}$.  Applying
\eqref{thm5_1} to $h_\eps$, with
$A=\sqrt{\sigma^2\Jfour(F)}$, gives
\[
 |\E[h(F)]-\E[h(G)]|
 \leq2\eps+2K_qA\eps^{-2}.
\]
The conclusion is immediate when $A=0$.  Otherwise, the choice
$\eps=(K_qA)^{1/3}$ yields
$|\E[h(F)]-\E[h(G)]|\leq4K_q^{1/3}A^{1/3}$.  Taking the supremum
over all one-Lipschitz $h$ proves \eqref{thm5_2} with
$L_q=4K_q^{1/3}$.
\qedhere

\end{proof}

\end{document}
